\SourceLeaf{364}%
simply as \(240:1823\). He would have had to do without this advantage if he had started from the fraction \(\frac{362}{209}\).

At the close of the memoir, Mr. M. also remarks that \textsc{Columella}, \textit{de re rustica} V, 2, made use of one of the approximate values of \(\BareRoot{3}\), namely \(\frac{26}{15}\). For the area of an equilateral triangle he takes the sum of one third and one tenth of the square described on its side.

